\documentclass[10pt,a4paper]{amsart}
\usepackage[margin=19mm]{geometry}
\usepackage{amsmath,amssymb,mathtools}
\usepackage[scr=rsfs,cal=euler]{mathalfa}
\usepackage{xfrac}
\usepackage[unicode,hidelinks,bookmarks=false]{hyperref}
\newcommand{\ie}{i.e.\ }
\newcommand{\cf}{cf.\ }
\newcommand{\resp}{respectively\ }
\newcommand{\Subsection}{\subsection}
\providecommand{\nobreakdash}{\nobreak\hbox{-}}
\setlength{\emergencystretch}{2em}
\multlinegap0pt
\usepackage{bm}
\usepackage{bbm}
\usepackage{dsfont}
\usepackage{xspace}
\usepackage{cancel}
\usepackage{mathtools}
\usepackage{amsthm}
\makeatletter
\@ifundefined{maintheorem}{\newtheorem{maintheorem}{Theorem}}{}
\@ifundefined{theorem}{\newtheorem{theorem}{Theorem}}{}
\@ifundefined{corollary}{\newtheorem{corollary}[theorem]{Corollary}}{}
\@ifundefined{lemma}{\newtheorem{lemma}[theorem]{Lemma}}{}
\makeatother
\DeclareMathOperator{\Aut}{Aut}
\DeclareMathOperator{\Epi}{Epi}
\DeclareMathOperator{\Ext}{Ext}
\DeclareMathOperator{\Homeo}{Homeo}
\DeclareMathOperator{\cd}{cd}
\DeclareMathOperator{\inc}{inc}
\newcommand{\Z}{\mathbb Z}
\renewcommand{\leq}{\leqslant}
\title[Aspherical 4-manifolds with elementary amenable fundamental ]{Aspherical 4-manifolds with elementary amenable fundamental group}
\author{James F. Davis, J. A. Hillman}
\date{Source: arXiv:2501.12512v2 (CC BY 4.0)}
\begin{document}
\maketitle

\noindent\emph{Source and licence.} Aspherical 4-manifolds with elementary amenable fundamental group. Version arXiv:2501.12512v2 (CC BY 4.0), distributed under the Creative Commons Attribution 4.0 International licence, as recorded in the arXiv licence metadata for this exact version and shown on the abstract page https://arxiv.org/abs/2501.12512v2. Full rights notice: licenses/arxiv-2501.12512-CC-BY-4.0.txt.\par\smallskip

\noindent\textit{Editorial scope.} Counted unit: the Theorem whose source label is \texttt{thmD} in the article (source0.tex lines 1210--1241, source label \texttt{thmD}), delivered together with the article's own complete original proof of it (source0.tex lines 1243--1313, one \texttt{proof} environment of 71 lines). No mathematics is written, repaired or re-derived: statement and proof are the article's own text line for line. Required context retained uncounted: A (lines 119--142); components of boundary (lines 368--377); cd<n (lines 391--419); FJC for EA in dim 4 (lines 717--719); epiaut (lines 1178--1182). Every definition, display and label of those lines is retained unchanged. Omitted from this package and not redistributed: 1--118, 143--367, 378--390, 420--716, 720--1177, 1183--1209, 1242--1242, 1314--1502. Nothing inside a retained excerpt is omitted or altered: every retained body is delivered exactly as the article's lines are written, so no omission receipt of any kind accompanies this package. Citations: the retained excerpts carry no citation command at all, so no bibliography entry of the article is transcribed and no reference is redistributed. Reference adaptation: none was needed and none was made; every reference inside the delivered unit resolves inside it, so no unresolved reference or citation is produced. Printed slips kept literal: no printed word, symbol, hypothesis or number of the counted statement or of its proof was altered. Layout and class: the article's own class is \texttt{amsart} with packages including amssymb, amscd, amsmath, amsthm, epsfig, latexsym, enumerate; the wrapper is the lane's pinned \texttt{amsart} editorial wrapper at 10pt on A4, which restates the theorem-like environments the retained text needs and adds the article's own macro definitions that the retained text needs (\texttt{\textbackslash Aut}, \texttt{\textbackslash Epi}, \texttt{\textbackslash Ext}, \texttt{\textbackslash Homeo}, \texttt{\textbackslash Z}, \texttt{\textbackslash cd}, \texttt{\textbackslash inc}, \texttt{\textbackslash leq}), with extra wrapper packages (bm, bbm, dsfont, xspace, cancel, mathtools). The delivered numbering is the unit's own. Assets: no retained span contains a figure environment, a graphics-inclusion command, a PDF-page inclusion, a table or a TikZ picture; no image file is copied into this package, the package declares no asset dependency and no image file is redistributed with it. Licence: version arXiv:2501.12512v2 of this article is distributed under the Creative Commons Attribution 4.0 International licence, as recorded in the arXiv licence metadata for this exact version and shown on the abstract page https://arxiv.org/abs/2501.12512v2. A full rights notice is delivered at licenses/arxiv-2501.12512-CC-BY-4.0.txt. Characters: every retained excerpt is pure ASCII, so the delivered engine, XeLaTeX with its default Latin Modern text font, reports no missing character.
 \begin{maintheorem} \label{A}
Let $M$ be a compact aspherical $4$-manifold with boundary $N$,
and let $\inc:N\to{M}$ be the inclusion.
If $\pi=\pi_1M$ is elementary amenable, then one of the following conditions holds:
\begin{enumerate}
 \setcounter{enumi}{-1}
\item$\pi=1$ and $N$ is a homology $3$-sphere;
\item$\pi\cong\mathbb{Z}$, $N$ is connected, and $\pi_1\inc$ 
is an epimorphism with perfect kernel 
(i.e.,  $N$ is a $\Z[\Z]$-homology $S^2\times{S^1}$);
\item$\pi\cong{BS(1,m)}$ for some $m\not=0$, $N$ is connected, 
$\pi_1\inc$ is an epimorphism and $H^i(\inc;\Z\pi)$ is an isomorphism
for $i\leq2$; 
\item$\pi$ is a polycyclic $PD_3$-group and either 
\begin{enumerate}
\item$N$ has 
two components and $\pi_1\inc_i$ is an epimorphism for $i=1,2$;
or 
\item$N$ is connected and $[\pi:\mathrm{Im}(\pi_1\inc)]=2$;
\end{enumerate}
 in each case the kernels are perfect; 
\item$\pi$ is a polycyclic $PD_4$-group and $N$ is empty.
\end{enumerate}
\end{maintheorem}

\begin{lemma} \label{components of boundary}
Let $M$ be a compact aspherical $n$-manifold with fundamental group $\pi$.
\label{boundary components}
\begin{enumerate}
\item \label{cdpin} $\cd \pi =n \Longleftrightarrow \partial M$ is empty.
\item \label{cdpin-2}$\cd \pi \leq{n-2} \Longrightarrow \partial M$ has one component.
\item \label{pdn-1} $ \pi$ is $PD_{n-1} \Longrightarrow \partial M$ has one or two components.
\qed
\end{enumerate}
\end{lemma}

\begin{theorem}
\label{cd<n}
Let $M$ be a compact aspherical $n$-manifold.
Then $\pi=\pi_1M$ is of type $F$ and $\cd\pi\leq4$.
\begin{enumerate}
\item If $\cd \pi < n$ and $\pi \not = 1$, then the image of $\pi_1\inc_i$ is infinite for all $i$.
\item If $\cd \pi \leq{n-2}$, then $\pi_1\inc$ is an epimorphism.
\item If $\cd \pi =n- 2$, then $\pi_1\inc$ has nontrivial kernel.
\item If $\cd \pi \leq1$, then  $\ker(\pi_1\inc)$ is a perfect group. 
\item If $\pi$ is a $PD_2$-group, then $\ker(\pi_1\inc)$ 
has abelianization  $\Z$.
\item If $\pi$ is a $PD_{n-1}$-group with orientation character $w_{\pi}$,
then each $\ker(\pi_1\inc_i)$ is a perfect group and either
\begin{enumerate}
\item $\partial M$ has two components and each $\pi_1\inc_i$ is an epimorphism and
$$
w_1M = w_\pi : \pi \to \Z^{\times}
$$ 
or
\item $\partial M$ has one component, 
$[\pi:\mathrm{Im}(\pi_1\inc)]=2$ and
$$
w_1M = w_{\pi} \cdot w_{\partial} : \pi \to \Z^{\times}
$$
where $w_\partial: \pi \to \Z^{\times}$ has kernel $\mathrm{Im}(\pi_1\inc)$.
\qed
\end{enumerate}
\end{enumerate}
\end{theorem}

\begin{corollary} \label{FJC for EA in dim 4}
If $\pi$ is the fundamental group of a compact aspherical 4-manifold and $\pi$ is elementary amenable, then the Farrell-Jones Conjecture in $K$- and $L$-theory holds.  In particular, the Borel Uniqueness Conjecture for such manifolds hold: a homotopy equivalence between compact aspherical 4-manifolds with elementary amenable fundamental group which is a homeomorphism on the boundary is homotopic, relative to the boundary, to a homeomorphism.
\end{corollary}

\begin{lemma}
\label{epiaut}
Let $\alpha,\beta:\nu\to\pi$ be epimorphisms with the same kernel $\kappa$.
Then there is an automorphism $\theta$ of $\pi$ such that $\theta\circ\alpha=\beta$.
\end{lemma}

\begin{theorem}
\label{thmD}
Let $\pi$ be a finitely presented elementary amenable group and $N$ a closed $3$-manifold.
Assume that $p:N\to{K(\pi,1)}$ is a map such that $\pi$, $N$ and $p$ 
satisfy one of the conditions of Theorem \ref{A} (with $p$ for $inc$ in case $(2)$).
Then 
\begin{enumerate}
\item{} compact aspherical $4$-manifolds $M$ with $\pi_1M\cong\pi$, $\partial{M}\cong{N}$
and $w_1M=w$ are classified by $\pi$ and $N$ alone if 
$\pi=1$ or $\Z$.
\end{enumerate}
In the other cases further invariants are needed:

\begin{enumerate}
\addtocounter{enumi}{1}
\item{}If $\pi\cong{BS(1,m)}$, $\nu=\pi_1N$, and
\begin{enumerate}
\item($m=\pm1$) then  the quotient of $\Epi(\nu,\pi)$  by the actions of $\Homeo(N,*)$ and $\Aut(\pi)$is trivial;
\item($|m|>1$) then there are at most two possibilities, corresponding to certain elements of the quotient of $\Epi(\nu,\pi)$ by the action of $\Homeo(N,*)$.  
\end{enumerate}
\item{If} $\pi$ is a polycyclic $PD_3$-group and
\begin{enumerate}
\item($N=N_1\sqcup{N_2}$, with $\nu_1=\pi_1N_1$ and $\nu_2=\pi_1N_2$)
then also an element of the quotient of 
$\Epi(\nu_1,\pi)\times{\Epi(\nu_2,\pi)}$ by the action of $\Homeo(N_1)\times{\Homeo(N_2)}$ and the diagonal action of $\Aut(\pi)$; 
\item($N$ connected, with $\nu=\pi_1N$) then also an element of the quotient of\\
 $\bigsqcup_{\{w\}}\Epi(\nu,\ker(w))$ 
by the actions of $\Homeo(N,*)$ and $\Aut(\pi)$, where the union is taken over all  epimorphisms $w:\pi\to\Z^\times$.
\end{enumerate}
\end{enumerate}
Each set of invariants satisfying the conditions of Theorem \ref{A} can be realized.
\end{theorem}

\begin{proof}
In cases (1) and (2), Lemma \ref{components of boundary} and Theorem \ref{cd<n}, show that $N$ is connected and $
\pi_1\inc$ is an epimorphism.   It then follows from Corollary \ref{FJC for EA in dim 4} that $M$ is determined by $\pi$, $N$ and the image of $\pi_1\inc$ in the quotient
of $\Epi(\pi_1N,\pi)$  by the actions of $\Homeo(N,*)$ and $\Aut(\pi)$.

If $\pi = 1$, then $\Epi(\pi_1N,\pi)$ is trivial, and if $\pi = \Z$, then Theorem \ref{A} shows that the quotient of $\Epi(\pi_1N,\pi)$ by $\Aut(\pi)$ is trivial.


%% 23 March 2026

We next need to explain the more explicit result in case (2b).
Suppose that $\pi=\pi_1M\cong{BS(1,m)}=\langle{a,t}\mid{tat^{-1}=a^m}\rangle$,  where $|m|>1$.
Then the generator $t$ determines an isomorphism $\pi/I(\pi)\cong\Z$,
and hence an isomorphism $\Z[\pi/I(\pi)]\cong\Lambda=\Z[t,t^{-1}]$.

Since $[\pi,\pi]$ is torsion-free and abelian, 
every homomorphism from $\nu$ to $\pi$ must factor through $\nu/I([\nu,\nu])$. 
The exact sequence of $(M,N)$ and Poincar\'e duality together imply that $\beta_1(N)=1$,
since $H_2(M,N)\cong{H^2(M;\Z^w)}$ is finite.
Hence $\nu/I(\nu)\cong\pi/I(\pi)$.
We assume now that we have fixed such an isomorphism.
The subgroup $A=I(\pi)\cong\Lambda/(t-m)\Lambda$ is torsion-free and of rank 1 
as an abelian group.
The exact sequence of $(M,N)$ with coefficients in $\Lambda=\Z[\pi/I(\pi)]$ 
gives a short exact sequence
\[
0\to{H_2(M,N;\Lambda)}\to{I(\nu)^{ab}}\to{A=I(\pi)}\to0,
\]
since $H_2(M;\Lambda)=H_2(A;\mathbb{Z})=A\wedge{A}=0$.
Now 
\[
H_2(M,N;\Lambda)\cong{H^2(M;{}^w\Lambda)}\cong
{\Ext^1_\Lambda(A,{}^w\Lambda)}\cong\Lambda/(mt-w(t))\Lambda,
\]
by Poincar\'e duality and the Universal Coefficient spectral sequence,
and so $I(\nu)^{ab}$ is torsion-free and of rank 2 as an abelian group.
Let $E$ be the preimage of $H_2(M,N;\Lambda)$ in $\nu$.
Then $E$ is characteristic in $\nu$, since
its image in $I(\nu)^{ab}$ is the kernel of multiplication by $mt-w(t)$.

If we use coefficients $\Z[\pi^{ab}]$ instead of $\Lambda$,
a similar argument shows that $[\nu,\nu]^{ab}$ has rank 2.
Hence the quotient of $[I(\nu),I(\nu)]$ by $I(\nu,\nu)$ is a torsion group.
Since $\pi$ is torsion-free and metabelian, the image of this quotient in $\pi$ is trivial.
Hence every epimorphism from $\nu$ to $\pi$ factors through $\nu/[I(\nu),I(\nu)]$.
In all cases, such epimorphisms have kernel $E$ and so induce isomorphisms 
$\nu/E\cong\pi$.
If $(M',N)$ is another such pair and the epimorphisms $p=\pi_1\inc$ 
and $p'=\pi_1\inc'$
induce the same isomorphism from $\nu/I(\nu)$ to $\pi/I(\pi)=\Z$ 
then $\ker(p)=\ker(p')$.
Hence we may again invoke Lemma \ref{epiaut}.

The effect of changing our choice of identification of $\nu/I(\nu)$ with $\pi/I(\pi)$
is to invert the action of $t$ on $I(\nu)^{ab}$, and thus to replace
$E$ by the preimage in $\nu$ of the kernel of multiplication by $t-w(t)m$.
Thus for such $\pi$ and $N$ there are two possible kernels.
(These may be equivalent under composition with automorphisms of $\nu$ 
induced by self-homeomorphisms of $N$.)


%% end 23 March 2026


In case 3(a) the group $\pi$ is a canonical quotient of each of $\nu_1$
and $\nu_2$, since $\pi$ is solvable and the kernels are perfect.
Hence the actions of $\Homeo(N_1,*)$ and $\Homeo(N_2,*)$ induce homomorphisms 
from these groups to $\Aut(\pi)$. 
If one of these homomorphisms is onto then $\Homeo(N_1,*)\times{\Homeo(N_2,*)}$ and   $\Aut(\pi)$ together act transitively on $\Epi(\nu_1,\pi)\times{\Epi(\nu_2,\pi)}$,
and the classification needs only $\pi$ and $N$.  
\end{proof}

\end{document}
